%% ma: L12-B01-0013
%% lop: 12
%% chuong: 1
%% bai: 1
%% dang: 12.1.8
%% loai: TN4
%% muc: TH
%% dap_an: "D"
%% nguon: "(NBV)-ĐỀ SỐ 9-MỖI NGÀY 1 ĐỀ THI 2025.pdf (D09, MỖI NGÀY 1 ĐỀ - Đề số 9), Phần I Câu 5"
%% kien_thuc: "Bài 1 (xét tính đơn điệu của hàm số bằng đạo hàm)"
%% trang_thai: chinh-thuc
%% ghi_chu: "Đáp án gốc D đúng (sympy: y' = 6(x+3)(x-4))."
\begin{ex}
Cho hàm số $y=2x^3-3x^2-72x-6$. Khẳng định nào sau đây đúng?
\choice{Hàm số nghịch biến trên khoảng $(-\infty;4)$}{Hàm số nghịch biến trên khoảng $(-3;+\infty)$}{Hàm số đồng biến trên khoảng $(-3;4)$}{\True Hàm số đồng biến trên khoảng $(4;+\infty)$}
\loigiai{$y'=6x^2-6x-72=6(x+3)(x-4)$; $y'=0\Leftrightarrow x=-3$ hoặc $x=4$. Hàm số đồng biến trên $(-\infty;-3)$ và $(4;+\infty)$, nghịch biến trên $(-3;4)$. Vậy chỉ có phương án D đúng.}
\end{ex}
